\section{ラグランジュ（Lagrange）の微分方程式}
次の形
\begin{gather}
  y = x f\Bigl(\frac{dy}{dx}\Bigr) + g\Bigl(\frac{dy}{dx}\Bigr)
\end{gather}
で表される一階非線形常微分方程式をラグランジュの微分方程式、
あるいはダランベールの微分方程式という。

特に$f(x) = x$のとき、この微分方程式は
\begin{gather}
  y = x \frac{dy}{dx} + g\Bigl(\frac{dy}{dx}\Bigr)
\end{gather}
と表されて、この微分方程式を特にクレロー（Clairaut）の微分方程式という。

クレローの微分方程式は両辺をさらに$x$で微分して
\begin{gather}
  \frac{dy}{dx} = \frac{dy}{dx} + x \frac{d^2y}{dx^2} + \frac{dg(y')}{dx} \\
  0 = xy'' + \frac{dg(y')}{dy'} y'' \\
  y''\left(x + \frac{dg(y')}{dy'}\right) = 0
\end{gather}
ゆえに、
\begin{gather}
  y'' = 0 \quad\text{または}\quad x + \frac{dg(y')}{dy'} = 0
\end{gather}
として解く。
前者は単に任意定数$C$を用いて
\begin{gather}
  \frac{dy}{dx} = C
\end{gather}
とできるので与えられた微分方程式より、
\begin{gather}
  y = Cx + g(C)
\end{gather}
として解を得る。これは1つの任意定数を持つので一般解である。

一方で後者は
\begin{gather}
  \begin{dcases*}
    x + \frac{dg(y')}{dy'} = 0 \\
    y = xy' + g(y')
  \end{dcases*}
\end{gather}
から連立させて$y'$を消去すれば任意定数を含まない解$y$を得る。
この解を特異解という。

ラグランジュの微分方程式も同様にまず両辺を$x$で微分して
\begin{gather}
  y' = f(y') + xy''\frac{df(y')}{dy'} + \frac{dg(y')}{dy'}y''
\end{gather}
ここで、$y' = p$とおけばこの微分方程式は
\begin{gather}
  p = f(p) + xp' f'(p) + p'g'(p) \\
  p = f(p) + \frac{dp}{dx} \bigl(xf'(p) + g'(p)\bigr) \\
  p - f(p) = \frac{dp}{dx} \bigl(xf'(p) + g'(p)\bigr) \\
  \frac{dx}{dp} = \frac{xf'(p)}{p - f(p)} + \frac{g'(p)}{p - f(p)} \\
  \frac{dx}{dp} + \frac{f'(p)}{f(p) - p} x + \frac{g'(p)}{f(p) - p} = 0
\end{gather}
と変形できるので、これは一階線形微分方程式に帰着する。
このとき、$f(p) - p \not\equiv 0$を使っているので、% [?] 一瞬でも0になる場合は？
関数$f$が$f(x) = x$となるような条件では
異なるもの（クレローの微分方程式）として異なる解法をとることも自然である。
